\setcounter{ex}{0}
\setcounter{paragraph}{0}
\PhanI
%\loai{Thời điểm ban đầu vật ở biên}
\begin{ex}
Một vật dao động điều hòa dọc theo trục $Ox$ với biên độ $5\ cm$, tần số góc $2\pi\ rad/s$. Pha ban đầu của dao động là $\pi\ rad$. Phương trình dao động của vật là
\choice
{\True $x=5\cos(2\pi t+\pi)\ cm$}
{$x=5\cos(2\pi t)\ cm$}
{$x=5\cos(\pi t+\pi)\ cm$}
{$x=5\cos(\pi t)\ cm$}
\loigiai{
}
\end{ex}
\begin{ex} %[3P1B4-2][Loai1] 
Một vật dao động điều hoà với biên độ $A = 2\ cm$, tần số $f = 2\ Hz$. Chọn gốc toạ độ O trùng với vị trí cân bằng của vật, chọn gốc thời gian $t = 0$ là lúc vật đạt li độ cực tiểu. Phương trình dao động của chất điểm là
\choice 
{$x = 2\cos\left(2t + \pi\right)\ cm$}
{$x = 2\cos\left(4\pi t\right)\ cm$}
{$x = 2\cos\left(2\pi t - \pi\right)\ cm$}
{\True $x = 2\cos\left(4\pi t - \pi\right)\ cm$}
\loigiai{
\begin{itemize}
\item Ta có: $\omega = 2.2\pi = 4\pi$ rad/s.
\item Gốc thời gian $t = 0$ là lúc vật đạt li độ cực đại về phía âm $\Rightarrow \varphi = \pi\Rightarrow$ Phương trình dao động của chẩt điểm là:
$x = 2\cos\left(4\pi t+ \pi\right)$ cm hay $x = 2\cos\left(4\pi t- \pi\right)$ cm.
\end{itemize}
}
%\haidong{5}
\end{ex} \haidong{15}

%\loai{Thời điểm đầu vật qua li độ $\dfrac{A}{2}$ hoặc $-\dfrac{A}{2}$} 
\begin{ex}%[3P1B4-2][Loai1]
Một vật dao động điều hòa dọc theo trục Ox trên quỹ đạo dài 10 cm. Vật thực hiện 90 dao động toàn phần trong 3 phút. Tại thời điểm $t = 0$ s vật đi qua vị trí có li độ 2,5 cm theo chiều dương. Phương trình dao động của vật là
\choice
{$x=10\cos\left(\pi t+\dfrac{2\pi}{3}\right)$ cm}
{$x=5\cos\left(2\pi t+\dfrac{\pi}{3}\right)$ cm}
{$x=5\cos\left(2\pi t-\dfrac{\pi}{3}\right)$ cm}
{\True $x=5\cos\left(\pi t-\dfrac{\pi}{3}\right)$ cm}
\loigiai{
\begin{itemize}
\item Ta  có: $T=\dfrac{\Delta t}{N}=2\ s\Rightarrow \omega=\dfrac{2\pi}{T}=\pi\ rad/s$.
\item Biên độ: $A=\dfrac{L}{2}=5\ cm$.
\item Tại thời điểm $t = 0$ s, $x_0= 2,5\ cm=\dfrac{A}{2}$ và $v_0>0$ nên $\varphi=-\dfrac{\pi}{3}\ rad$.
\end{itemize}
}
%\haidong{6}
\end{ex} \haidong{15}
\begin{ex}
Một vật dao động điều hòa với biên độ $A = 4\ cm$ và chu kì $T = 2\ s$, chọn gốc thời gian là lúc vật đi qua vị trí cân bằng theo chiều âm. Phương trình dao động của vật là
\choice
{\True $x = 4 \cos \left(\pi t + \dfrac{\pi}{2}\right)\ cm$}
{$x = 4 \cos \left(2 \pi t - \dfrac{\pi}{2}\right)\ cm$}
{$x = 4 \cos \left(\pi t - \dfrac{\pi}{2}\right)\ cm$}
{$x = 4 \cos \left(2 \pi t + \dfrac{\pi}{2}\right)\ cm$}
\loigiai{}
\end{ex}

%\loai{Thời điểm đầu vật qua li độ $\dfrac{A\sqrt3}{2}$ hoặc $-\dfrac{A\sqrt3}{2}$} 
\begin{ex}%[3P1B4-2]
Một vật dao động điều hòa dọc theo trục Ox với biên độ 6 cm, tần số 2 Hz. Tại thời điểm $t = 0$ s vật đi qua vị trí li độ $-3\sqrt{3}$ cm và đang chuyển động lại gần vị trí cân bằng. Phương trình dao động của vật là
\choice
{\True $x=6\cos\left(4\pi t-\dfrac{5\pi}{6}\right)$ cm}
{$x=6\cos\left(4\pi t-\dfrac{2\pi}{3}\right)$ cm}
{$x=6\cos\left(4\pi t+\dfrac{5\pi}{6}\right)$ cm}
{$x=6\cos\left(4\pi t-\dfrac{\pi}{6}\right)$ cm}
\loigiai{
\begin{itemize}
\item Biên độ: $A = 6$ cm.
\item Tần số góc $\omega =2\pi f=4\pi \ rad/s$. 
\item Tại $t = 0$: $x = -3\sqrt{3}\ cm =  \dfrac{-A\sqrt{3}}{2}$  theo chiều dương $\Rightarrow \phi  =  -\dfrac{5\pi}{6}$. 
\end{itemize}
}
%\haidong{7}
\end{ex} \haidong{15}
%\loai{Biết li độ và vận tốc tại thời điểm $t=0$. Viết phương trình} 
\begin{ex} %[3P1B4-2][Loai1] 
Một chất điểm dao động điều hòa trên trục Ox. Trong thời gian 31,4 s chất điểm thực hiện được 50 dao động toàn phần. Gốc thời gian là lúc chất điểm đi qua vị trí có li độ 3 cm theo chiều âm với tốc độ là 10$\sqrt{3}$ cm/s. Lấy $\pi = 3,14$. Phương trình dao động của chất điểm là
\choice 
{\True $x = 2\sqrt{3}\cos\left(10t + \dfrac{\pi}{6}\right)$ cm}
{$x = 3\sqrt{2}\cos\left(20t + \dfrac{\pi}{3}\right)$ cm}
{$x = 3\sqrt{2}\cos\left(10t - \dfrac{\pi}{3}\right)$ cm}
{$x = 2\sqrt{3}\cos\left(20t + \dfrac{\pi}{6}\right)$ cm}
\loigiai{
\begin{itemize}
\item Ta có: $T = \dfrac{31,4}{50} = \dfrac{\pi}{5}\Rightarrow \omega = 10\ rad/s\Rightarrow x = A\cos\left(10 t+ \varphi\right)\ cm\Rightarrow v = -10.A\sin\left(10 t+ \varphi\right)\ cm/s$.
\item Tại $t = 0$, $x = 3$ cm và $v = -10\sqrt{3}$ cm/s, nên ta có hệ phương trình:
$\begin{cases}
3 = A\cos\varphi\quad (1)
-10\sqrt{3} = -10.A\sin\varphi\quad (2)
\end{cases}$
\item Lấy $(2)$ chia $(1)$ theo vế, được: $\tan\varphi = \dfrac{1}{\sqrt{3}}\Rightarrow \varphi = \dfrac{\pi}{6}$ hoặc $ \varphi = \dfrac{7\pi}{6}$.
\item Do vật chuyển động theo chiều âm nên chọn $\varphi = \dfrac{\pi}{6}$.
\item Thay vào $(1)$, được $3 = A\cos\dfrac{\pi}{6} \Rightarrow A = 2\sqrt{3}\ cm \Rightarrow x = 2\sqrt{3}\cos\left(10t+ \dfrac{\pi}{6}\right)$ cm.
\end{itemize}
}
%\haidong{5}
\end{ex} \haidong{15}
%\loai{Biết gia tốc và vận tốc tại thời điểm $t=0$. Viết phương trình x} 
\begin{ex} %[3P1K4-2][Loai1] 
Một chất điểm dao động điều hòa với chu kì bằng 1 s. Gốc tọa độ ở vị trí cân bằng. Tại thời điểm ban đầu ($t = 0$), chất điểm có vận tốc bằng $4\pi\sqrt{3}$ cm/s và gia tốc bằng $8\pi^2\ cm/s^2$. Phương trình dao động của chất điểm là
\choice 
{\True $x = 4\cos\left(2\pi t - \dfrac{2\pi}{3}\right)$ cm}
{$x = 2\cos\left(4\pi t + \dfrac{\pi}{3}\right)$ cm}
{$x = 4\cos\left(2\pi t + \dfrac{2\pi}{3}\right)$ cm}
{$x = 2\cos\left(4\pi t - \dfrac{2\pi}{3}\right)$ cm}
\loigiai{
\begin{itemize}
\item Ta có: $\omega = \dfrac{2\pi}{1} = 2\pi\ rad/s\Rightarrow x = A\cos\left(2\pi t+ \varphi\right)\ cm
\Rightarrow v = -2\pi.A\sin\left(2\pi t+ \varphi\right)\ cm/s
\Rightarrow a = -4\pi^2.A\cos\left(2\pi t+ \varphi\right)\ cm/s$.
\item Tại $t = 0$:
$\begin{cases}
4\pi \sqrt{3} = -2\pi A \sin\varphi\quad (1)\\
8\pi^2 = -4.\pi^2 A \cos\varphi\quad (2)
\end{cases}$
\item Lấy $(1)$ chia $(2)$ theo vế, ta được:
$\tan\varphi = \sqrt{3}\Rightarrow \varphi = \dfrac{\pi}{3}$ hoặc $\varphi = \dfrac{4\pi}{3}$.
\item Mà tại $t = 0$, $v > 0$ chọn nên $\varphi = \dfrac{4\pi}{3}$.
\item Thay giá trị trên vào (1), tính được $A = 4\ cm\Rightarrow x = 4\cos\left(2\pi t+\dfrac{4\pi}{3}\right)\ cm\Rightarrow x = 4\cos\left(2\pi t-\dfrac{2\pi}{3}\right)$ cm.
\end{itemize}
}
%\haidong{6}
\end{ex} \haidong{15}
%\loai{Viết phương trình vận tốc}
\begin{ex} %[3P1B4-2][Loai1] 
Một chất điểm dao động điều hòa trên trục Ox. Trong thời gian 31,4 s chất điểm thực hiện được 100 dao động toàn phần. Gốc thời gian là lúc chất điểm đi qua vị trí có li độ 2 cm theo chiều âm với tốc độ là 40$\sqrt{3}$ cm/s. Lấy $\pi = 3,14$. Phương trình vận tốc của chất điểm là
\choice 
{$v = 40\cos\left(10t + \dfrac{2\pi}{3}\right)$ cm/s}
{$v = 40\cos\left(20t + \dfrac{5\pi}{6}\right)$ cm/s}
{$v = 80\cos\left(10t + \dfrac{2\pi}{3}\right)$ cm/s}
{\True $v = 80\cos\left(20t + \dfrac{5\pi}{6}\right)$ cm/s}
\loigiai{
\begin{itemize}
\item Ta có: $T = \dfrac{31,4}{100} = \dfrac{\pi}{10}\ s \Rightarrow \omega = 20\ rad/s\Rightarrow x = A\cos\left(20t+ \varphi\right)\ cm\Rightarrow v = -20.A\sin\left(20t+ \varphi\right)\ cm/s$.
\item Tại $t = 0$, $x = 2$ cm và $v = -40\sqrt{3}$ cm/s, nên ta có hệ phương trình:
$\begin{cases}
2 = A\cos\varphi\quad (1)\\
-40\sqrt{3} = -20.A\sin\varphi\quad (2)
\end{cases}$
\item Lấy $(2)$ chia $(1)$ theo vế, được: $\tan\varphi = \sqrt{3}\Rightarrow \varphi = \dfrac{\pi}{3}$ hoặc $ \varphi = \dfrac{4\pi}{3}$.
\item Do vật chuyển động theo chiều âm nên chọn $\varphi = \dfrac{\pi}{3}$. 
\item Thay vào (1), được $2 = A\cos\dfrac{\pi}{3} \Rightarrow A = 4\ cm
\Rightarrow v = -80\sin\left(20t+ \dfrac{\pi}{3}\right)\ cm
\Rightarrow v = 80\cos\left(20t+ \dfrac{5\pi}{6}\right)\ cm/s$.
\end{itemize}
}
%\haidong{10}
\end{ex} \haidong{15}


%\loai{Loại khác}
\begin{ex}%[3P1K4-2][Loai1]
Một chất điểm đang dao động điều hòa trên trục Ox với tần số bằng 0,5 Hz và biên độ bằng 2 cm quanh vị trí cân bằng là gốc tọa độ O. Tại thời điểm ban đầu, $t = 0$, chất điểm có li độ dương và đang chuyển động với vận tốc $\pi\sqrt{2}$ cm/s. Phương trình dao động của chất điểm là
\choice 
{$x = 2\cos\left(\pi t - \dfrac{2\pi}{3}\right)$ cm}
{$x = 2\cos\left(\pi t - \dfrac{3\pi}{4}\right)$ cm}
{$x = 2\cos\left(\pi t + \dfrac{2\pi}{3}\right)$ cm}
{\True $x = 2\cos\left(\pi t - \dfrac{\pi}{4}\right)$ cm}
\loigiai{
\begin{itemize}
\item Ta có: $\omega = 0,5.2\pi = \pi\ rad/s\Rightarrow x = 2\cos\left(\pi t+ \varphi\right)\ cm
\Rightarrow v = -2\pi\sin\left(\pi t+ \varphi\right)\ cm/s$.
\item Tại $t = 0$, $\pi\sqrt{2} = -2\pi.\sin\varphi$ $\Rightarrow \sin\varphi = \dfrac{-1}{\sqrt{2}} = \sin\left(-\dfrac{\pi}{4}\right)$
$\Rightarrow \varphi = -\dfrac{\pi}{4}$ hoặc $\varphi = \dfrac{5\pi}{4}$.
\item Mà tại $t = 0$, $x > 0 $ nên $\varphi = -\dfrac{\pi}{4}
\Rightarrow x = 2\cos\left(\pi t- \dfrac{\pi}{4}\right)$ cm.
\end{itemize}
}
%\haidong{10}
\end{ex} \haidong{15}
\begin{ex}
Đồ thị dao động điều hòa của một vật như hình vẽ. Phương trình dao động của vật là 
\immini{\choice
{\True $x = 8 \cos \left(5 \pi t - \dfrac{\pi}{2}\right)\ (cm)$}
{$x = 8 \cos \left(10 \pi t + \dfrac{\pi}{2}\right)\ (cm)$}
{$x = 8 \cos \left(5 \pi t + \dfrac{\pi}{2}\right)\ (cm)$}
{$x = 8 \cos \left(10 \pi t - \dfrac{\pi}{2}\right)\ (cm)$}
}{
\begin{hinhanh}{7}
\begin{tikzpicture}[very thick,>=stealth']
\pgfplotsset{width=7cm,height=4.5cm,compat=1.4}
\begin{axis}[
axis lines=middle, xmin=0, xmax=1.7*pi, xstep=1,
ymin=-4.5, ymax=5, ystep=0.1,
grid=major,%Vẽ chế độ lưới theo trục tọa độ Oxy,
x grid style={dashed, black},
axis line style={very thick,Mapcolor}, % <--- dòng thêm vào
xlabel=$t~(s)$,
xlabel style={above},
xtick ={1,2,3,4,5,6,7,8},
xticklabels={},
y grid style={dashed, black},
extra y ticks={0},
extra y tick label={0},
ylabel=$x~(cm)$,
ylabel style={above},
ytick={-4,0, 4},
yticklabels={,,8},]
\addplot[line width=1.2pt, domain=0:1.5*pi, samples=400, Mapcolor]{4*cos(deg(0.5*pi*x-pi/2))};
\addplot[Mapcolor,mark=*,mark options={fill=white,mark size=1.5pt}, nodes near coords={0,3}]
coordinates {(3, -4)};
\end{axis}
\end{tikzpicture}\end{hinhanh}}
\loigiai{ }
\end{ex}
\begin{ex}%[3P1-19.1K][Loai1]
\immini[thm]{Cho một chất điểm dao động điều hòa, sự phụ thuộc của li độ vào thời gian được biểu diễn trên đồ thị như hình vẽ. 
Biên độ và pha ban đầu của dao động lần lượt là

\choice
{8 cm; $\dfrac{\pi}{2}$ rad}
{8 cm; $-\dfrac{\pi}{2}$ rad}
{\True 4 cm; $\dfrac{\pi}{2}$ rad}
{4 cm; $-\dfrac{\pi}{2}$ rad}}{\begin{tikzpicture}[draw=Mapcolor,scale=1.15,thick,>=stealth']
\draw[->] (0,0)--(7,0)node[below]{$t\ (s)$};
\draw[->] (0,-2.3)--(0,2.3)node[right]{$x\ (cm)$};
\draw[dashed] (0,1.5)--(6,1.5)--(6,0)(0,-1.5)--(6,-1.5)--(6,0);
\draw [Mapcolor, line width = 1.2pt, domain=0:6, samples=500]%
plot (\x, {1.5*cos(\x*120+90)});
\node at (0,0)[left]{$O$};
\node at (0,1.5)[left]{$4$};
\node at (0,-1.5)[left]{$-4$};
\node at (1.5,0)[below]{$1$};
\node at (3,0)[below]{$2$};
\node at (4.5,0)[below]{$3$};
\end{tikzpicture}}
\loigiai{
\begin{itemize}
\item Từ đồ thị ta suy ra: $A=4$ cm.
\item Tại $t = 0$ vật qua vị trí cân bằng theo chiều âm $\Rightarrow$ pha ban đầu của dao động là $\dfrac{\pi}{2}$.
\end{itemize}
}
%\haidong{3}
\end{ex} \haidong{20}
\begin{ex}%[3P1-19.1K][Loai1]
\immini[thm]{Cho một chất điểm dao động điều hòa, sự phụ thuộc của li độ vào thời gian được biểu diễn trên đồ thị như hình vẽ. Phương trình li độ của chất điểm là

\choice
{$x=6\cos\left(3\pi t + \dfrac{\pi}{3}\right)\ cm$}
{$x=6\cos\left(2,5\pi t + \dfrac{\pi}{3}\right)\ cm$}
{$x=6\cos\left(3\pi t - \dfrac{\pi}{3}\right)\ cm$}
{\True $x=6\cos\left(2,5\pi t - \dfrac{\pi}{3}\right)\ cm$}}{\begin{tikzpicture}[draw=Mapcolor,scale=1.15,thick,>=stealth']
\draw[->] (0,0)--(6.5,0)node[below]{$t\ (s)$};
\draw[->] (0,-2.3)--(0,2.3)node[right]{$x\ (cm)$};
\draw[dashed] (0,1.5)--(5.66,1.5)--(5.66,0)(0,-1.5)--(5.66,-1.5)--(5.66,0);
\draw [Mapcolor, line width = 1.2pt, domain=0:5.666, samples=500]%
plot (\x, {1.5*cos(\x*90-60)});
\node at (0,0)[left]{$O$};
\node at (0,1.5)[left]{$6$};
\node at (0,0.75)[left]{$3$};
\node at (0,-1.5)[left]{$-6$};
\node at (1.666,0)[below]{$\dfrac{1}{3}$};
\end{tikzpicture}}
\loigiai{
\begin{itemize}
\item Từ đồ thị $\Rightarrow$ $ A = 6 $ cm.
\item Thời điểm ban đầu $ t = 0 $, chất điểm có li độ $3\ cm = \dfrac{A}{2}$ và đang đi về biên dương\\ $\Rightarrow$ pha ban đầu $\varphi = -\dfrac{\pi}{3}$ rad.
\item Từ thời điểm ban đầu tới thời điểm $t=\dfrac{1}{3}\ s$, chất điểm đi từ vị trí 3 cm theo chiều dương qua biên dương và tới vị trí cân bằng theo chiều âm nên nó quay được một góc $\Delta \varphi = \dfrac{\pi}{3} + \dfrac{\pi}{2} = \dfrac{5\pi}{6}\Rightarrow \Delta t = \dfrac{5T}{12} = \dfrac{1}{3} \Rightarrow T = 0,8\ s \Rightarrow \omega = \dfrac{2\pi}{T} = 2,5\pi\ rad/s$.
\item Phương trình dao động của chất điểm là $x = 6\cos\left(2,5\pi t - \dfrac{\pi}{3}\right)\ cm$.
\end{itemize}
}
%\haidong{4}
\end{ex} \haidong{20}
\begin{ex}%[3P1-19.1K][Loai1]
\immini[thm]{Cho một chất điểm dao động điều hòa, sự phụ thuộc của li độ vào thời gian được biểu diễn trên đồ thị như hình vẽ. 
Phương trình vận tốc của chất điểm là
\choice
{$v=32\pi \cos\left(4\pi t + \pi \right)\ cm/s$}
{$v=32\pi \cos\left(4\pi t\right)\ cm/s$}
{$v=64\pi \cos\left(8\pi t\right)\ cm/s$}
{\True $v=64\pi \cos\left(8\pi t + \pi \right)\ cm/s$}}{\begin{tikzpicture}[draw=Mapcolor,scale=1.15,thick,>=stealth']
\draw[->] (0,0)--(7,0)node[below]{$t\ (s)$};
\draw[->] (0,-2.3)--(0,2.3)node[right]{$x\ (cm)$};
\draw[dashed] (0,1.5)--(6,1.5)--(6,0)(0,-1.5)--(6,-1.5)--(6,0);
\draw [Mapcolor, line width = 1.2pt, domain=0:6, samples=500]%
plot (\x, {1.5*cos(\x*60+90)});
\node at (0,0)[left]{$O$};
\node at (0,1.5)[left]{$8$};
\node at (0,-1.5)[left]{$-8$};
\node at (6,0)[below left]{$0,25$};
\end{tikzpicture}}
\loigiai{
\begin{itemize}
\item Từ đồ thị ta thấy: $ A = 8 $ cm; $ T = 0,25 $ s.
\item Tại $ t = 0 $ thì vật đang ở vị trí cân bằng và chuyền động theo chiều âm.
\item Phương trình li độ của vật: $x=8\cos\left(8\pi t + \dfrac{\pi} {2}\right)\ cm$\\
$\Rightarrow$ Phương trình vận tốc của vật: $v=64\pi \cos\left(8\pi t + \pi \right)\ cm/s$.
\end{itemize}
}
%\haidong{8}
\end{ex} \haidong{20}
\begin{ex}%[3P1-19.1K][Loai1] 
\immini[thm]{
Một vật dao động điều hòa trên trục Ox có đồ thị biểu diễn sự phụ thuộc của li độ x vào thời gian t như hình vẽ, pha ban đầu của dao động là
\choice
{$10\pi t-\dfrac{\pi }{2}$}
{$10\pi t+\dfrac{\pi }{2}$}
{$-\dfrac{\pi }{2}$}
{\True $+\dfrac{\pi }{2}$}
}{
\pgfplotsset{width=7.5cm,height=4cm,compat=1.4}
\begin{tikzpicture}[draw=Mapcolor,scale=1,thick,>=stealth']
\begin{axis}[
axis lines=middle,
xmin=0,xmax=4.5,xstep=1,ymin=-2.16,ymax=2.5,ystep=0.1,
grid=major, %Có các tùy chọn: minor|major|both|none
x grid style={dashed},
y grid style={dashed},
axis line style={very thick,Mapcolor},
ytick={-2,0,2},
xtick={1,2,3,4,5,6},
xticklabels={,{0,2} },
yticklabels={},
xticklabel style={Mapcolor,below},
xlabel style={above=3pt},
xlabel=$t(s)$,
ylabel style={right=3pt},
ylabel={$x$}
]
\addplot[line width=1.2pt,domain=0:4,samples=200,Mapcolor]{2*cos(deg(0.5*pi*x+pi/2))};
\end{axis}
\end{tikzpicture}
}
\loigiai{
Ban đầu vật đi qua vị trí cân bằng theo chiều âm $\varphi_0=0,5\pi$.
}
%\haidong{3}
\end{ex} \haidong{20}
\begin{ex}%[3P1-19.1K][Loai1]
\immini[thm]{Cho một chất điểm dao động điều hòa, sự phụ thuộc của li độ vào thời gian được biểu diễn trên đồ thị như hình vẽ. 
Phương trình li độ của chất điểm là
\choice
{$x=3\cos(4\pi t - \dfrac{\pi}{3})\ cm$}
{$x=3\cos(2\pi t + \dfrac{\pi}{3})\ cm$}
{\True $x=3\cos(2\pi t - \dfrac{\pi}{3})\ cm$}
{$x=3\cos(4\pi t + \dfrac{\pi}{6})\ cm$}}{\begin{tikzpicture}[draw=Mapcolor,scale=1,thick,>=stealth']
\draw[->] (0,0)--(6.5,0)node[below]{$t\ (s)$};
\draw[->] (0,-2.3)--(0,2.3)node[right]{$x\ (cm)$};
\draw[dashed] (0,1.5)--(5.66,1.5)--(5.66,0)(0,-1.5)--(5.66,-1.5)--(5.66,0)(0.66,1.5)--(0.66,0);
\draw [Mapcolor, line width = 1.2pt, domain=0:5.666, samples=500]%
plot (\x, {1.5*cos(\x*90-60)});
\node at (0,0)[left]{$O$};
\node at (0,1.5)[left]{$3$};
\node at (0,0.75)[left]{$1,5$};
\node at (0,-1.5)[left]{$-3$};
\node at (0.666,0)[below]{$\dfrac{1}{6}$};
\end{tikzpicture}}
\loigiai{
\begin{itemize}
\item Từ đồ thị ta thấy: $A=3\ cm;\  \dfrac{T}{6}=\dfrac{1}{6}\ s \Rightarrow T=1\ s$.
\item Tại $ t = 0 $ thì vật đang đi qua vị trí $x=\dfrac{A}{2}$ theo chiều dương
$\Rightarrow$ Phương trình li độ của vật: $x=3\cos\left(2\pi t - \dfrac{\pi}{3}\right)\ cm$.
\end{itemize}
}
%\haidong{3}
\end{ex} \haidong{20}
\begin{ex}%[3P1-19.1K][Loai1]
\immini[thm]{Hình vẽ bên là đồ thị biểu diễn sự phụ thuộc của li độ x vào thời gian t của một vật dao động điều hòa.
Tốc độ cực đại của vật là
\choice
{5,24 cm/s}
{1,05 cm/s}
{\True 3,14 cm/s}
{6,28 cm/s}}{
\pgfplotsset{width=7.5cm,height=4cm,compat=1.4}
\begin{tikzpicture}[draw=Mapcolor,scale=1,thick,>=stealth']
\begin{axis}[
axis lines=middle,
xmin=0,xmax=4.5,xstep=1,ymin=-2.16,ymax=2.5,ystep=0.1,
grid=major, %Có các tùy chọn: minor|major|both|none
x grid style={dashed},
y grid style={dashed},
ytick={-2,0,1,2},
xtick={1,2,3,4,5,6},
xticklabels={1,2},
yticklabels={$-2$,0,$+1$,$+2$},
xticklabel style={black,below},
xlabel style={above=3pt},
xlabel=$t(s)$,
ylabel style={right=3pt},
ylabel={$x(cm)$}
]
\addplot[line width=1.2pt,domain=0:4,samples=200,Mapcolor]{2*cos(deg(0.5*pi*x-pi/3))};
\end{axis}
\end{tikzpicture}
}
\loigiai{
\begin{itemize}
\item Từ đồ thị, ta thấy tại $t = 0$, vật đi qua vị trí $x = 1$ cm theo chiều dương.
\item Tại thời điểm $t = 0,5$ s, vật đi qua vị trí cân bằng theo chiều âm.
\item Từ hình vẽ, ta xác định được $T=4\ s \Rightarrow \omega =\dfrac{\pi}{2}$ rad/s.
\item Tốc độ cực đại của vật $v_{\max}=\omega A=3,14$ cm/s.
\end{itemize}
}
%\haidong{4}
\end{ex} \haidong{20}
\begin{ex}%[3P1-19.1K][Loai1]
\immini[thm]{Cho một chất điểm dao động điều hòa, sự phụ thuộc của vận tốc vào thời gian được biểu diễn trên đồ thị như hình vẽ. 
Phương trình li độ của chất điểm là
\choice
{\True $x=4\cos\left(2\pi t - \dfrac{\pi} {2}\right)\ cm$}
{$x=4\cos\left(2\pi t + \dfrac{\pi} {2}\right)\ cm$}
{$x=8\cos\left(\pi t - \dfrac{\pi} {2}\right)\ cm$}
{$x=8\cos\left(\pi t + \dfrac{\pi} {2}\right)\ cm$}}{\begin{tikzpicture}[draw=Mapcolor,scale=1,thick,>=stealth']
\draw[->] (0,0)--(7,0)node[below]{$t\ (s)$};
\draw[->] (0,-2.3)--(0,2.3)node[right]{$v\ (cm/s)$};
\draw[dashed] (0,1.5)--(6,1.5)--(6,0)(0,-1.5)--(6,-1.5)--(6,0);
\draw [Mapcolor, line width = 1.2pt, domain=0:6, samples=500]%
plot (\x, {1.5*cos(\x*120)});
\node at (0,0)[left]{$O$};
\node at (0,1.5)[left]{$8\pi$};
\node at (0,-1.5)[left]{$-8\pi$};
\node at (0.75,0)[below]{$0,25$};
\end{tikzpicture}}
\loigiai{
\begin{itemize}
\item Từ đồ thị ta thấy: $\dfrac{T}{4}=0,25\ s\Rightarrow T=1$ s $\Rightarrow$ $\omega = 2\pi\ rad/s$.
\item Tại $ t = 0 $ thì vận tốc của vật đang có "giá trị" cực đại ở biên dương\\ 
$\Rightarrow$ x có pha ban đầu $\varphi = -\dfrac{\pi}{2}$ ( do x chậm $\dfrac{\pi}{2}$ so với v).
\item $A = \dfrac{v_{max}}{\omega} = 4\ cm$
$\Rightarrow$ Phương trình li độ: $x=4\cos\left(2\pi t - \dfrac{\pi} {2}\right)\ cm$.
\end{itemize}
}
%\haidong{3}
\end{ex} \haidong{20}
\begin{ex}%[3P1-19.1K][Loai1]
\immini[thm]{Sự phụ thuộc của li độ theo thời gian của một chất điểm dao động điều hòa được biểu diễn như hình vẽ. 
Tại điểm nào trong các điểm M, N, K, H có vectơ gia tốc và vectơ vận tốc ngược chiều nhau?

\choice
{Điểm M}
{\True Điểm K}
{Điểm N}
{Điểm H}}{\begin{tikzpicture}[draw=Mapcolor,scale=1,thick,>=stealth']
\draw[->] (0,0)--(7,0)node[below]{$t\ (s)$};
\draw[->] (0,-2.3)--(0,2.3)node[right]{$x\ (cm)$};
\draw[dashed] (0,1.5)--(6,1.5)--(6,0)(0,-1.5)--(6,-1.5)--(6,0);
\draw [Mapcolor, line width = 1.2pt, domain=0:6, samples=500]%
plot (\x, {1.5*cos(\x*120)});
\draw(0.75,0)circle(1pt)node[above right]{M};
\draw(1.5,-1.5)circle(1pt)node[below]{N};
\draw(2.5,0.75)circle(1pt)node[above left]{K};
\draw(3.5,0.75)circle(1pt)node[above right]{H};
\end{tikzpicture}}
\loigiai{
Từ đồ thị ta xét các vị trí:
\begin{itemize}
\item Tại M vật đang đi qua vị trí cân bằng theo chiều âm $\Rightarrow$ gia tốc của vật có độ lớn $ = 0$.
\item Tại N vật đang ở vị trí biên âm $\Rightarrow$ vận tốc của vật có độ lớn $=0$.
\item Tại K vật đang có li độ x và đang chuyển động từ vị trí cân bằng ra biên dương\\ $\Rightarrow$ Tại K vectơ gia tốc hướng ngược chiều dương, vectơ vận tốc hướng cùng chiều dương
\item Tại H vật đang có li độ x và đang chuyển động từ vị trí biên dương về vị trí cân bằng \\
$\Rightarrow$ Tại K vectơ gia tốc hướng ngược chiều dương, vectơ vận tốc hướng ngược chiều dương.
\end{itemize}
}
%\haidong{3}
\end{ex} \haidong{20}
\setcounter{ex}{0}
\PhanII
\begin{ex}
Một vật dao động điều hòa trên trục $Ox$. Hình bên là đồ thị biểu diễn sự phụ thuộc của li độ $x$ vào thời gian $t$.
\immini{\choiceTF[t]
{\True Chu kì dao động của vật bằng $2\ s$}
{Pha ban đầu của vật bằng $0$}
{\True Phương trình dao động điều hòa của vật là $x = 10 \cos(\pi t + \pi)\ cm$}
{Vận tốc của vật tại thời điểm $t = 0,5\ s$ là $-10\pi\ m/s$}}{
\begin{hinhanh}{7}
\begin{tikzpicture}[draw=Mapcolor,very thick,>=stealth']
\pgfplotsset{width=7cm,height=4.5cm,compat=1.4}
\begin{axis}[
axis lines=middle, xmin=0, xmax=2*pi, xstep=1,
ymin=-4.5, ymax=5, ystep=0.1,
grid=major,%Vẽ chế độ lưới theo trục tọa độ Oxy,
x grid style={dashed, Mapcolor},
axis line style={very thick,Mapcolor}, % <--- dòng thêm vào
xlabel=$t~(s)$,
xlabel style={above},
xtick ={1,2,3,4,5,6,7,8},
xticklabels={},
y grid style={dashed, Mapcolor},
extra y ticks={0},
extra y tick label={0},
ylabel=$x~(cm)$,
ylabel style={above},
ytick={-4,0, 4},
yticklabels={$-10$,,10},]
\addplot[line width=1.2pt, domain=0:1.7*pi, samples=400, Mapcolor]{4*cos(deg(0.5*pi*x-pi))};
\addplot[Mapcolor,mark=*,mark options={fill=white,mark size=1.5pt}, nodes near coords={0,2}]
coordinates {(1, -4)};
\addplot[Mapcolor,mark=*,mark options={fill=white,mark size=1.5pt}, nodes near coords={1,0}]
coordinates {(2, -4)};
\addplot[Mapcolor,mark=*,mark options={fill=white,mark size=1.5pt}, nodes near coords={1,5}]
coordinates {(3, -4)};
\addplot[Mapcolor,mark=*,mark options={fill=white,mark size=1.5pt}, nodes near coords={2,0}]
coordinates {(4, -4)};
\end{axis}
\end{tikzpicture}
\end{hinhanh}}
\loigiai{
\begin{itemchoice}
\itemch	
\itemch	
\itemch	
\itemch	
\end{itemchoice}
}
\end{ex}
\begin{ex}
Cho phương trình dao động điều hòa: $x=4 \cos \left(4 \pi t - \dfrac{\pi}{6}\right)\ (cm)$.
\choiceTF[t]
{\True Chiều dài quỹ đạo là $8\ cm$}
{\True Chu kì của dao động bằng $0,5\ s$}
{Tốc độ cực tiểu bằng $16 \pi\ cm/s$}
{\True Khi pha dao động bằng $\dfrac{2\pi}{3}$ thì gia tốc bằng $32\pi^{2}\ cm/s^2$}
\loigiai{
\begin{itemchoice}
\itemch	
\itemch	
\itemch	
\itemch	
\end{itemchoice}
}
\end{ex}
\begin{ex}
Một vật dao động điều hòa trên trục Ox. Hình bên là đồ thị biểu diễn sự phụ thuộc của li độ $x$ vào thời gian $t$.
\immini{\choiceTF[t]
{\True Biên độ của dao động là $4\ cm$}
{Vận tốc của vật lúc $t=2,5\ s$ có giá trị là $-\sqrt{2}\ cm/s$}
{\True Phương trình dao động của vật có dạng $x=4\cos\left(\dfrac{\pi}{2}t+\pi\right)\ (cm;\ s)$}
{Chu kì dao động của vật là $1\ s$}}
{\begin{tikzpicture}[draw=Mapcolor,very thick,>=stealth']
\pgfplotsset{width=7cm,height=4.5cm,compat=1.4}
\begin{axis}[
axis lines=middle, xmin=0, xmax=2*pi, xstep=1,
ymin=-4.5, ymax=5, ystep=0.1,
grid=major,%Vẽ chế độ lưới theo trục tọa độ Oxy,
x grid style={dashed, Mapcolor},
axis line style={very thick,Mapcolor}, % <--- dòng thêm vào
xlabel=$t~(s)$,
xlabel style={above},
xtick ={1,2,3,4,5,6,7,8},
xticklabels={},
y grid style={dashed, Mapcolor},
extra y ticks={0},
extra y tick label={0},
ylabel=$x~(cm)$,
ylabel style={above},
ytick={-4,0, 4},
yticklabels={$-4$,,4},]
\addplot[line width=1.2pt, domain=0:1.7*pi, samples=400, Mapcolor]{4*cos(deg(0.5*pi*x))};
\addplot[Mapcolor,mark=*,mark options={fill=white,mark size=1.5pt}, nodes near coords={1}]
coordinates {(1, -4)};
\addplot[Mapcolor,mark=*,mark options={fill=white,mark size=1.5pt}, nodes near coords={2}]
coordinates {(2, -4)};
\addplot[Mapcolor,mark=*,mark options={fill=white,mark size=1.5pt}, nodes near coords={3}]
coordinates {(3, -4)};
\addplot[Mapcolor,mark=*,mark options={fill=white,mark size=1.5pt}, nodes near coords={4}]
coordinates {(4, -4)};
\addplot[Mapcolor,mark=*,mark options={fill=white,mark size=1.5pt}, nodes near coords={5}]
coordinates {(5, -4)};
\end{axis}
\end{tikzpicture}}

\loigiai{
\begin{itemchoice}
\itemch	
\itemch	
\itemch	
\itemch	
\end{itemchoice}
}
\end{ex}
\begin{ex}
Cho vật dao động điều hoà có đồ thị li độ - thời gian như hình vẽ.
\immini{\choiceTF[t]
{\True Dao động của vật có biên độ: $A = 10\ cm$ và tần số góc: $\omega = \dfrac{\pi}{2}\ rad/s$}
{\True Gia tốc của vật tại thời điểm $3,5\ s$ có giá trị: $a = -12,5\sqrt{2}\ cm/s^{2}$}
{Vận tốc của vật tại thời điểm $3,5\ s$ có giá trị $-5\sqrt{5}\ cm/s$}
{\True Phương trình dao động: $x = 10 \cos\left(\dfrac{\pi}{2} t\right)\ (cm)$}
}{
\begin{hinhanh}{7}\begin{tikzpicture}[draw=Mapcolor,very thick,>=stealth']
\pgfplotsset{width=7cm,height=4.5cm,compat=1.4}
\begin{axis}[
axis lines=middle, xmin=0, xmax=2*pi, xstep=1,
ymin=-4.5, ymax=5, ystep=0.1,
grid=major,%Vẽ chế độ lưới theo trục tọa độ Oxy,
x grid style={dashed, Mapcolor},
axis line style={very thick,Mapcolor}, % <--- dòng thêm vào
xlabel=$t~(s)$,
xlabel style={above},
xtick ={1,2,3,4,5,6,7,8},
xticklabels={},
y grid style={dashed, Mapcolor},
extra y ticks={0},
extra y tick label={0},
ylabel=$x~(cm)$,
ylabel style={above},
ytick={-4,0, 4},
yticklabels={$-10$,,10},]
\addplot[line width=1.2pt, domain=0:1.7*pi, samples=400, Mapcolor]{4*cos(deg(0.5*pi*x))};
\addplot[Mapcolor,mark=*,mark options={fill=white,mark size=1.5pt}, nodes near coords={1}]
coordinates {(1, -4)};
\addplot[Mapcolor,mark=*,mark options={fill=white,mark size=1.5pt}, nodes near coords={2}]
coordinates {(2, -4)};
\addplot[Mapcolor,mark=*,mark options={fill=white,mark size=1.5pt}, nodes near coords={3}]
coordinates {(3, -4)};
\addplot[Mapcolor,mark=*,mark options={fill=white,mark size=1.5pt}, nodes near coords={4}]
coordinates {(4, -4)};
\addplot[Mapcolor,mark=*,mark options={fill=white,mark size=1.5pt}, nodes near coords={5}]
coordinates {(5, -4)};
\end{axis}
\end{tikzpicture}\end{hinhanh}}
\loigiai{
\begin{itemchoice}
\itemch	
\itemch	
\itemch	
\itemch	
\end{itemchoice}
}
\end{ex}















\setcounter{ex}{0}
\PhanIII
\begin{ex}
Một vật chuyển động tròn đều trên đường tròn đường kính 12 cm. Hình chiếu của vật lên một đường kính dao động điều hoà với biên độ là bao nhiêu cm?
\shortans[oly]{6}
\loigiai{}
\end{ex}
\begin{ex}
Một vật dao động điều hòa trong nửa chu kì đi được quãng đường $10\ cm$. Khi vật có li độ $3\ cm$ thì có vận tốc $16\pi\ cm/s$. Chu kì dao động của vật là bao nhiêu giây (làm tròn kết quả đến chữ số hàng phần mười)?
\shortans[oly]{0,5}
\loigiai{
Từ dữ kiện bài toán, tính được $A=10\ cm$, dùng $v=\omega\sqrt{A^2-x^2}$ để tính $\omega$, từ đó suy ra $T\approx 0,5\ s$.
}
\end{ex}
\begin{ex}%book%[3P1B4-1][Loai1]
Một chất điểm dao động điều hoà theo hàm cosin với chu kì 2 s và có vận tốc $- 1$ m/s vào lúc pha dao động bằng $\dfrac{\pi}{4}$ rad. Biên độ dao động là bao nhiêu cm?
\shortans[oly]{45}
\loigiai{
\begin{itemize}
\item Ta có: $\omega  = \pi$  rad/s.
\item Khi $\phi =\dfrac{\pi}{4}:\ x =\dfrac{A\sqrt{2}}{2}\ (-) \Rightarrow v=-\dfrac{v_{1\max}\sqrt{2}}{2}=-\dfrac{\omega A\sqrt{2}}{2}=-1\Rightarrow A=45\ cm$.
\end{itemize}
}
%\haidong{4}
\end{ex} \haidong{15}
\begin{ex}%book%[3P1KC-1][Loai1]
Một con lắc dao động điều hòa theo phương ngang với chu kì T và biên độ 10 cm. Biết ở thời điểm t vật có li độ 6 cm, ở thời điểm $t + \dfrac{T}{2}$ vật có tốc độ 80 cm/s. Tần số góc của dao động là bao nhiêu rad/s?
\shortans[oly]{10}
\loigiai{
\begin{itemize}
\item $\Delta t=\dfrac{T}{2}$: ngược pha $\Rightarrow$  $x_2=-x_1=-6$ cm.
\item Lại có: $x _{2}^2+\dfrac{v_2^2}{{{\omega}^2}}=A^2\Rightarrow \omega =10$ rad/s.
\end{itemize}
}
%\haidong{5}
\end{ex} \haidong{20}
\begin{ex}%book%[3P1-12.2K][Loai1]
Một chất điểm dao động điều hoà theo phương ngang. Biết ở thời điểm t vật có tốc độ 40 cm/s, sau đó ba phần tư chu kì gia tốc của vật có độ lớn $1,6\pi$ $m/s^2$. Tần số dao động của vật bằng bao nhiêu Hz?
\shortans[oly]{2}
\loigiai{
\begin{itemize}
\item Ta có: $\Delta t=\dfrac{3T}{4}$: vuông pha $\Rightarrow\left| a_2 \right|=\omega \left|v_1\right|\Rightarrow \omega =4\pi \Rightarrow f = 2\ Hz$.
\end{itemize}

}
%\haidong{7}
\end{ex} \haidong{20}
\begin{ex}%book%[3P1KC-2][Loai1]
Một chất điểm dao động điều hòa theo phương ngang với chu kì T và biên độ $5\ cm$. Biết ở thời điểm t vật có tốc độ $20\ cm/s$, ở thời điểm $t +\dfrac{T}{4}$ gia tốc của vật có độ lớn $1\ m/s^2$. Li độ tại thời điểm t có độ lớn bằng bao nhiêu cm?
\shortans[oly]{3}
\loigiai{
\begin{itemize}
\item Ta có: $\Delta t=\dfrac{T}{4}$: vuông pha $\Rightarrow \left| a_2 \right|=\omega \left| v_1 \right|\Rightarrow \omega =5$ rad/s.
\item Lại có: $x _1^2+\dfrac{v_1^2}{\omega^2}=A^2\Rightarrow \left| x_1 \right|=3$ cm.
\end{itemize}
}
%\haidong{9}
\end{ex} \haidong{20}






























